\section{Gabarito e resoluções comentadas — Fiscalização de Contratos de TI}

\begin{solucao}{15.01}\textbf{Gabarito: A.} A IN SGD/ME nº 94/2022 organiza o macroprocesso em Planejamento da Contratação, Seleção do Fornecedor e Gestão do Contrato. ETP e TR são artefatos/etapas dentro do planejamento, não fases equivalentes ao macroprocesso.\end{solucao}
\begin{solucao}{15.02}\textbf{Gabarito: A.} O ETP parte da necessidade, compara alternativas e sustenta a viabilidade da solução. Não substitui contrato, recebimento ou aplicação de sanções.\end{solucao}
\begin{solucao}{15.03}\textbf{Gabarito: A.} O TR transforma a solução escolhida em requisitos executáveis e mensuráveis, incluindo modelo de execução/gestão e pagamento. Marca específica sem justificativa pode restringir competição.\end{solucao}
\begin{solucao}{15.04}\textbf{Gabarito: A.} SLA define níveis de serviço com métricas, metas, janelas e consequências conforme contrato. Não é sistema de licitação nem licença.\end{solucao}
\begin{solucao}{15.05}\textbf{Gabarito: A.} Pesquisa de preços sustenta estimativa razoável, documentando fontes, comparabilidade e método. Não escolhe fornecedor antecipadamente nem se resume ao menor número encontrado.\end{solucao}
\begin{solucao}{15.06}\textbf{Gabarito: A.} Recebimento definitivo deve decorrer de verificação de conformidade. Mera entrega física pode fundamentar recebimento provisório, conforme objeto/procedimento, mas não substitui aceite final.\end{solucao}
\begin{solucao}{15.07}\textbf{Gabarito: A.} Riscos mudam com solução, seleção, execução e mudança contratual; o mapa precisa ser vivo. Tratar risco só no início reduz valor do controle.\end{solucao}
\begin{solucao}{15.08}\textbf{Gabarito: A.} Lock-in é dependência de fornecedor, formato, API, licença ou tecnologia que torna saída cara/demorada. Não é bloqueio de conta nem incidente de rede.\end{solucao}
\begin{solucao}{15.09}\textbf{Gabarito: A.} Indicador auditável precisa ser reproduzível: fórmula, fonte, período, meta e regras de inclusão/exclusão. Opinião do fiscal sem critério é insuficiente.\end{solucao}
\begin{solucao}{15.10}\textbf{Gabarito: A.} Segregação reduz possibilidade de uma pessoa definir, executar e validar operação sem revisão. Não é burocracia como fim em si.\end{solucao}
\begin{solucao}{15.11}\textbf{Gabarito: A.} ETP deve avaliar alternativas antes de concluir pela solução. Decisão prévia seguida de justificativa pode mascarar direcionamento e inviabilizar análise de vantajosidade.\end{solucao}
\begin{solucao}{15.12}\textbf{Gabarito: A.} Medir apenas postos pode remunerar ociosidade e deslocar risco de produtividade ao órgão. O modelo deve relacionar pagamento a entregas/resultados quando tecnicamente possível.\end{solucao}
\begin{solucao}{15.13}\textbf{Gabarito: A.} Em 30 dias há 43.200 minutos. A indisponibilidade para 99,9\% é $0,1\%\times43.200=43,2$ min. A regra real depende da janela e exclusões definidas.\end{solucao}
\begin{solucao}{15.14}\textbf{Gabarito: A.} Sem definição de manutenção programada, o fornecedor pode excluir períodos relevantes e inflar a disponibilidade. Métrica precisa de regra objetiva anterior à medição.\end{solucao}
\begin{solucao}{15.15}\textbf{Gabarito: A.} A parte economicamente interessada controla a única evidência, o que fragiliza independência e verificabilidade. O órgão deve ter logs, acesso ou mecanismo de reconciliação.\end{solucao}
\begin{solucao}{15.16}\textbf{Gabarito: A.} Três propostas não garantem mercado representativo quando há vínculo entre cotantes. O auditor deve verificar independência, comparabilidade e possíveis sinais de conluio.\end{solucao}
\begin{solucao}{15.17}\textbf{Gabarito: A.} A IN 65 estabelece parâmetros e prioriza fontes públicas relevantes quando disponíveis. Ignorá-las sem justificativa pode distorcer a estimativa.\end{solucao}
\begin{solucao}{15.18}\textbf{Gabarito: A.} Plano de saída trata exportação, formatos, documentação, prazos, suporte à transição, apagamento e continuidade, reduzindo lock-in.\end{solucao}
\begin{solucao}{15.19}\textbf{Gabarito: A.} Aceite deve confrontar entrega com critérios previamente definidos. A declaração da própria contratada não é suficiente para provar conformidade.\end{solucao}
\begin{solucao}{15.20}\textbf{Gabarito: A.} RTO precisa ser traduzido em arquitetura, responsabilidades, procedimentos e testes. Cláusula abstrata sem capacidade demonstrável é frágil. RPO é objetivo distinto.\end{solucao}
\begin{solucao}{15.21}\textbf{Gabarito: A.} Uso de 24\% das licenças indica necessidade de revisar premissas, sazonalidade, crescimento e alternativas de licenciamento. O pagamento formal não prova economicidade.\end{solucao}
\begin{solucao}{15.22}\textbf{Gabarito: A.} Penalidade exige previsão, fato comprovado, nexo com obrigação e devido processo/contraditório aplicável. Automatismo sem evidência é inadequado.\end{solucao}
\begin{solucao}{15.23}\textbf{Gabarito: A.} Prorrogação deve ser motivada e demonstrar aderência às condições legais e vantajosidade, considerando desempenho e risco de transição. Não é direito automático.\end{solucao}
\begin{solucao}{15.24}\textbf{Gabarito: A.} Contrato deve disciplinar finalidade/instruções, acesso, segurança, incidentes, suboperadores, retenção, transferência e eliminação conforme papéis e bases aplicáveis.\end{solucao}
\begin{solucao}{15.25}\textbf{Gabarito: A.} OS/OFB precisa permitir rastrear demanda e verificar resultado. Autorizações vagas aumentam risco de escopo não controlado e pagamento indevido.\end{solucao}
\begin{solucao}{15.26}\textbf{Gabarito: A.} Receber logs brutos protegidos e recalcular amostras melhora independência da evidência. Isso não elimina todo risco, pois configuração e completude dos logs também precisam ser validadas.\end{solucao}
\begin{solucao}{15.27}\textbf{Gabarito: A.} Exceções sem taxonomia e evidência podem transformar qualquer falha em exclusão. O contrato deve definir força maior, manutenção, terceiros e processo de validação.\end{solucao}
\begin{solucao}{15.28}\textbf{Gabarito: A.} Exigir marca/modelo sem fundamento funcional/técnico pode restringir competição. O requisito deve descrever desempenho e compatibilidade ou justificar exceção.\end{solucao}
\begin{solucao}{15.29}\textbf{Gabarito: A.} Ponto de função mede tamanho funcional, não qualidade do código. Aceite deve combinar tamanho/entrega com defeitos, testes, segurança e critérios de qualidade.\end{solucao}
\begin{solucao}{15.30}\textbf{Gabarito: A.} Cloud tem componentes variáveis de custo. Ignorar egress, suporte, crescimento, logs e serviços associados subestima TCO e prejudica comparação de alternativas.\end{solucao}
\begin{solucao}{15.31}\textbf{Gabarito: A.} SLO extremo exige redundância e custo compatíveis. Meta arbitrária pode elevar preço ou ser inexequível. Planejamento deve relacionar necessidade, impacto e capacidade técnica.\end{solucao}
\begin{solucao}{15.32}\textbf{Gabarito: A.} E-mail genérico não demonstra quantidade, qualidade ou cumprimento de SLA. O fiscal precisa de evidência suficiente e apropriada ao objeto.\end{solucao}
\begin{solucao}{15.33}\textbf{Gabarito: A.} Concentrar especificação, seleção e aceite reduz checks and balances e aumenta risco de favorecimento/erro não detectado.\end{solucao}
\begin{solucao}{15.34}\textbf{Gabarito: A.} Risco é dinâmico. Mudança de arquitetura, fornecedor, legislação ou ameaça deve provocar reavaliação e atualização de respostas.\end{solucao}
\begin{solucao}{15.35}\textbf{Gabarito: A.} Código/documentação sem direitos e condições claras cria dependência do fornecedor e pode impedir manutenção por terceiros. O planejamento deve tratar propriedade/licenciamento e transição.\end{solucao}
\begin{solucao}{15.36}\textbf{Gabarito: A.} Mensalidade e licença perpétua têm horizontes e componentes diferentes; devem ser normalizadas por período, escopo, suporte e TCO. Média/mediana não corrige dados incomparáveis.\end{solucao}
\begin{solucao}{15.37}\textbf{Gabarito: A.} Métrica unilateral de fechamento pode ser manipulada. Reabertura, tempo até solução e aceite do usuário reduzem gaming e alinham indicador ao resultado.\end{solucao}
\begin{solucao}{15.38}\textbf{Gabarito: A.} Offboarding incompleto mantém acesso de terceiro sem necessidade. Deve haver revogação de contas, tokens, VPNs, certificados e privilégios no encerramento/transição.\end{solucao}
\begin{solucao}{15.39}\textbf{Gabarito: A.} Retenção demonstra existência potencial de cópias; só restore test evidencia que cadeia, chaves, catálogo e procedimentos realmente recuperam o serviço.\end{solucao}
\begin{solucao}{15.40}\textbf{Gabarito: A.} Emergência recorrente pode ser consequência de planejamento inadequado, dependência contratual ou falha de governança. O auditor deve examinar causa, não apenas preço da última contratação.\end{solucao}
\begin{solucao}{15.41}\textbf{Gabarito: A.} Sem portabilidade e APIs documentadas, o órgão perde poder de saída. Plano de exit deve estar no ETP/TR antes da contratação, não apenas quando o fornecedor aumenta custos.\end{solucao}
\begin{solucao}{15.42}\textbf{Gabarito: A.} Ausência de catálogo e SLA mensurável impede demonstrar o que foi comprado e entregue. O auditor deve cruzar demanda, pagamentos, chamados, equipe e resultados.\end{solucao}
\begin{solucao}{15.43}\textbf{Gabarito: A.} Evidência usada para pagar não deve poder ser reescrita unilateralmente. Telemetria protegida/imutável, sob domínio do órgão ou reconciliada, aumenta confiabilidade.\end{solucao}
\begin{solucao}{15.44}\textbf{Gabarito: A.} SOC deve ser contratado por capacidade e resultado de detecção/resposta, com cobertura de fontes, casos de uso, tempos, qualidade e evidências. Apenas quantidade de pessoas não mede segurança.\end{solucao}
\begin{solucao}{15.45}\textbf{Gabarito: A.} Ociosidade de cerca de 82\% sugere dimensionamento sem base ou antecipação excessiva. O achado deve examinar premissas, custo evitável e flexibilidade contratual.\end{solucao}
\begin{solucao}{15.46}\textbf{Gabarito: A.} Reequilíbrio depende do regime legal, matriz de risco e comprovação de evento/impacto. Risco ordinário do contratado não gera direito automático.\end{solucao}
\begin{solucao}{15.47}\textbf{Gabarito: A.} Uso secundário de dados pode violar finalidade, confidencialidade e direitos de propriedade, além de criar risco de memorização/exposição. Deve haver limites contratuais claros.\end{solucao}
\begin{solucao}{15.48}\textbf{Gabarito: A.} Encerramento precisa prever devolução/exportação, eliminação de produção e backups conforme retenção, evidência da destruição e exceções legais.\end{solucao}
\begin{solucao}{15.49}\textbf{Gabarito: A.} Resposta inicial pode ser cosmeticamente rápida enquanto usuário permanece indisponível. Métricas devem refletir restauração/resolução e severidade.\end{solucao}
\begin{solucao}{15.50}\textbf{Gabarito: A.} Fiscal não deve alterar critério para consumir orçamento. Pagamento deve corresponder ao objeto aceito; não conformidade precisa ser registrada e tratada conforme contrato.\end{solucao}
